\noindent\textbf{Elección del procedimiento.} La pérdida de altura libera energía potencial, mientras la resistencia realiza trabajo negativo. El balance es $mv_2^2/2-mv_1^2/2=mg\Delta z-W_D$, donde $W_D$ es aquí el módulo positivo de la pérdida. Multiplicar por $2/m$ deja el cuadrado de la velocidad.

\begin{enumerate}

\item Sumar la energía gravitatoria y restar la disipada con el signo indicado.

\[\fitmath{v2sq={v_{1}}^{2}+2\cdot g\cdot \mathrm{dz}-\frac{2\cdot \mathrm{WD}}{m}}\]

\[\fitmath{v2sq={(25)}^{2}+2\cdot (9.81)\cdot (600)-\frac{2\cdot (1500000)}{(400)}\approx 4897\,\mathrm{m^{2}/s^{2}}}\]

\item El radicando es positivo y se toma la raíz positiva porque se pide celeridad.

\[\fitmath{v_{2}=\sqrt{v2sq}}\]

\[\fitmath{v_{2}=\sqrt{(4897)}\approx 69.97857\,\mathrm{m/s}}\]

\end{enumerate}
